\section{Experimental Methodology}
\label{sec:methodology}

We evaluate the system through nine controlled experiments, each testing a
specific hypothesis about the governance model, agentic components, or
end-to-end behaviour.

\subsection{Hypotheses}

\begin{description}
  \item[H1] \kdd~governance increases the traceability score of generated
    artefacts compared to ungoverned generation.
  \item[H2] \sdd-constrained generation produces higher document compliance
    rates than free-form generation.
  \item[H3] \mcp-mediated tool access increases audit coverage compared to
    direct tool invocation.
  \item[H4] Skill-based agent invocation increases output consistency compared
    to free-prompt invocation.
  \item[H5] \rag/\cag~augmentation increases groundedness and citation
    coverage compared to no retrieval.
  \item[H6] Evidence-grounded copilot recommendations achieve higher quality
    scores than non-grounded recommendations.
  \item[H7] Digital twin validation blocks a measurable fraction of unsafe
    recommendations.
  \item[H8] The \rcc~dashboard achieves satisfactory usability scores in a
    simulated race engineering scenario.
\end{description}

\subsection{Experiment Design}

Each experiment follows a within-subjects design with two conditions: baseline
(without the proposed mechanism) and variant (with the mechanism). Metrics are
computed over synthetic datasets to enable full reproducibility without
confidential telemetry.

\subsection{Datasets}

All experiments use synthetic datasets generated by
\texttt{scripts/generate-synthetic-telemetry.py} (see
Appendix~\ref{app:experiment-configurations}). The datasets are documented in
\texttt{datasets/} with dataset cards specifying their structure, generation
method, and limitations.

\subsection{Statistical Analysis}

We report mean and standard deviation across ten independent runs per condition.
Statistical significance is assessed using paired Wilcoxon signed-rank tests at
$\alpha = 0.05$.
